% Stochastic PINNs: Feynman-Kac paths and path signatures.
% Target venues: Journal of Computational Physics, CMAME.
%
% This file exists from week one ON PURPOSE. Before running an experiment, write
% the sentence it is meant to support and the caption of the figure that will
% show it. If neither can be written, the experiment is not yet worth running.
%
% Build:  latexmk -pdf main.tex

\documentclass[11pt]{article}

\usepackage[margin=1in]{geometry}
\usepackage{amsmath,amssymb,amsthm}
\usepackage{graphicx}
\usepackage{booktabs}
\usepackage{subcaption}
\usepackage[numbers,sort&compress]{natbib}
\usepackage[colorlinks=true,allcolors=blue]{hyperref}
\usepackage{cleveref}

% figures/ is a symlink (or copy) of ../outputs/figures -- never hand-edited
\graphicspath{{figures/}}

\newcommand{\pinn}{PINN}
\newcommand{\fkpinn}{FK-PINN}
\newcommand{\E}{\mathbb{E}}
\newcommand{\R}{\mathbb{R}}
\newcommand{\dd}{\mathrm{d}}

\title{Stochastic Physics-Informed Neural Networks:\\
       Feynman--Kac Path Signatures for High-Dimensional PDEs}
\author{Federico Ganzer\\ \small Department of Engineering, University of Cambridge}
\date{\today}

\begin{document}
\maketitle

\begin{abstract}
% Write this LAST, but draft a placeholder now so the claim stays visible.
% One sentence each: problem, gap, method, headline number, implication.
% PLACEHOLDER -- written last, drafted now so the claim stays visible.
%
% One sentence each:
%   Problem     -- PINNs solve PDEs mesh-free but suffer stiff gradients and scale
%                  badly with dimension because the AD Laplacian costs d backward passes.
%   Gap         -- The Feynman-Kac representation replaces the pointwise residual with a
%                  path expectation, but the resulting estimator is noisy.
%   Method      -- We train on FK path expectations and summarise each path by its
%                  truncated signature rather than its endpoint.
%   Result      -- [NUMBER] relative $L^2$ at [DIMENSION] d, versus [NUMBER] for a
%                  standard PINN at equal wall-clock.
%   Implication -- [WHAT THIS UNLOCKS]

Placeholder abstract.

\end{abstract}

\section{Introduction}
\label{sec:intro}

% Structure to fill:
%
% 1. High-dimensional parabolic PDEs matter (finance, control/HJB, molecular dynamics)
%    and classical mesh-based methods die at d >~ 4.
% 2. PINNs are mesh-free but have two documented pathologies: stiff/imbalanced gradients
%    between loss terms, and an AD Laplacian costing d backward passes per collocation
%    point.
% 3. Feynman-Kac gives an alternative training signal -- an expectation over SDE paths --
%    that is dimension-robust by construction. Deep BSDE \citep{han2018deepbsde} is the
%    closest prior art; state precisely what is different here.
% 4. The FK estimator is noisy. Contribution: path signatures as a richer, lower-variance
%    summary of each path, plus the variance-reduction machinery to isolate that effect.
% 5. Contributions, as a bulleted list, each one pointing at a specific figure or table.

Placeholder.

\section{Background}
\label{sec:background}

% 2.1 PINNs. Loss = AD residual + BC/IC. Cite Raissi et al. 2019. State the two
%     pathologies you attack, with citations: gradient imbalance (Wang et al. 2021)
%     and the O(d) cost of the AD Laplacian.
% 2.2 Feynman-Kac. State the representation ONCE, with the exact sign and time
%     convention used throughout the paper, matching spinns/pdes/base.py. Every
%     later equation refers back to this one.
% 2.3 Path signatures. Iterated integrals, Chen's relation, truncation, log-signature,
%     and why time augmentation is mandatory for SDE paths. Cite Lyons; Chevyrev &
%     Kormilitzin 2016 for the ML-facing primer.
% 2.4 Related work. Deep BSDE (Han, Jentzen, E 2018), Walk-on-Spheres, neural CDEs/RDEs.
%     Be explicit and generous about what Deep BSDE already does -- the delta is the
%     contribution, and understating prior art is the fastest route to a hostile review.

Placeholder.

\section{Method}
\label{sec:method}

% 3.1 The FK training objective. Give the estimator, state whether the target is
%     detached, and prove or state its unbiasedness. Be explicit that
%     E[(u - Yhat)^2] = (u - E Y)^2 + Var(Yhat)/M, so the loss has a noise floor at
%     Var/M -- a reader who sees the loss plateau above zero must already know why.
% 3.2 Path features. Endpoint vs. full signature vs. log-signature. Give the feature
%     dimension explicitly (sum_k d^k) and confront its growth in d.
% 3.3 Variance reduction. Control variate with the network's own value as control;
%     antithetic pairing; importance sampling with ESS monitoring.
% 3.4 Algorithm box: one pseudocode listing of a single FK-PINN training step.
% 3.5 Cost model. Per-step cost of \pinn{} vs \fkpinn{} as a function of d, M and K.
%     This is what makes the equal-wall-clock comparison in Sec.~\ref{sec:experiments}
%     interpretable rather than merely empirical.

Placeholder.

\section{Experimental protocol}
\label{sec:experiments}

% Fix this section BEFORE running the comparisons. Everything here is decided in
% advance; that is what makes the results defensible.
%
% Baselines (all four, on every PDE):
%   1. analytic / high-accuracy reference (ground truth)
%   2. plain MC Feynman-Kac, no network (the accuracy ceiling \fkpinn{} should approach)
%   3. standard AD-based \pinn{} (the thing being beaten)
%   4. Deep BSDE (closest prior art)
%
% Controls: identical width/depth, optimiser, LR schedule, evaluation grid and seed set
% across methods. Only the loss/feature mechanism varies.
%
% Budget: report BOTH equal-gradient-steps and equal-wall-clock.
%
% Metrics: relative $L^2$ and $L^\infty$ on the fixed grid; error vs.\ wall-clock;
% error vs.\ dimension $d$; variance of the loss estimator per step; steps-to-target-error.
%
% Statistics: >= 5 seeds for every headline number; tables report mean $\pm$ std,
% plots show median with IQR band.
%
% Reproducibility: every number comes from a run directory recording the resolved
% config, git commit and seed; every figure is regenerated by scripts/make_figures.py.

Placeholder.

\section{Results}
\label{sec:results}

% Tables are \input from outputs/tables/*.tex -- generated by
% scripts/aggregate_results.py, never typed by hand, so numbers cannot drift
% away from the runs that produced them.
%
% \input{../outputs/tables/main_results.tex}
%
% Figures come from outputs/figures/ via scripts/make_figures.py:
%   loss_curves, error_vs_step, error_vs_time, estimator_variance, results_summary
%
% Write the caption BEFORE the experiment runs. If you cannot, the experiment is
% not yet worth running.

Placeholder.

\section{Discussion}
\label{sec:discussion}

% Say the honest things here, first, before a referee does:
%
% - Where the method LOSES. If \fkpinn{} is slower per step in low d, say so plainly
%   and show the crossing point in d. A paper that only wins is not believed.
% - Whether the signature effect survives once variance reduction is switched on.
%   If it does not, that is still a result and still publishable -- reported as
%   "signatures act as implicit variance reduction, and are subsumed by explicit
%   control variates" it is a genuine contribution.
% - Expect the signature effect to be SMALLER on the heat equation than on ADR:
%   with zero drift, every path is a Brownian bridge and there is less path shape
%   for a signature to summarise. State this prediction, then test it.
% - Signature dimension grows as sum_k d^k -- the honest scaling limitation.
% - Discretisation bias from finite Delta t, and the missed-excursion bias in
%   boundary-stopped paths.

Placeholder.

\section{Conclusion}
\label{sec:conclusion}

Placeholder.


\bibliographystyle{plainnat}
\bibliography{refs}

\end{document}
